\documentclass[11pt,a4paper]{article}
\usepackage[margin=1in]{geometry}
\usepackage{amsmath,amssymb,amsthm}
\usepackage{booktabs}
\usepackage{multirow}
\usepackage{graphicx}
\usepackage{microtype}
\emergencystretch=4em
\usepackage{hyperref}
\hypersetup{colorlinks=true,linkcolor=black,urlcolor=blue,citecolor=black}

\theoremstyle{plain}
\newtheorem{proposition}{Proposition}
\newtheorem{lemma}[proposition]{Lemma}
\theoremstyle{definition}
\newtheorem{remark}[proposition]{Remark}

\newcommand{\R}{\mathbb{R}}
\newcommand{\spec}{\operatorname{spec}}
\newcommand{\supp}{\operatorname{supp}}
\newcommand{\cond}{\kappa}
\newcommand{\blkdiag}{\operatorname{blkdiag}}

\title{From the Hessian to the Covariance:\\
Generalising the CBMR standard-error step}
\author{Strategies, proofs, speed-ups, and what it costs in memory}
\date{25 August 2026}

\begin{document}
\maketitle

\begin{abstract}
A companion note derives the CBMR observed information in closed form for all three of NiMARE's
observation distributions. That removes the Hessian as a bottleneck and immediately exposes the
next one: \texttt{CBMRModel.covariance} spends its remaining time in two dense LAPACK calls, a
singular value decomposition for a diagnostic condition number and an LU inverse. At the
21-group shape the SVD alone costs $163$~s against the $1.6$~s the Hessian now takes.

This note treats those two calls. We prove five propositions --- structural separation, the
blockwise inverse, the bordered (Schur) inverse, the block-diagonal spectrum, and the reduction
of the condition number to a symmetric eigenproblem --- and close each one symbolically with
\texttt{sympy}, exactly as the companion note does for the Hessians. We then establish which
designs each applies to and verify the resulting dispatch over the same nine-design,
three-distribution suite. One candidate is rejected on measurement rather
than on theory: LAPACK's \texttt{pocon} condition estimator is $18$--$264\times$ off the true
$2$-norm value here, comfortably inside its worst-case bound and comfortably too far to report
to a user. What generalises without any hypothesis on the design is simpler --- the information
matrix is symmetric, so a symmetric eigensolver replaces the SVD exactly.

We then measure memory, which the earlier work had not. On the real 21-group diagnosis pool
($8{,}420$ coefficients, $421$ bases, $21{,}789$ voxels, $8{,}178$ experiments) a complete run
--- load, fit, covariance, and twenty one-versus-rest contrasts --- peaks at $2.67$~GB and takes
$44$~s, against $5.49$~GB and $769$~s for the chunked-autodiff path and a $28.7$~GB intermediate
for stock NiMARE, which is why the model could not be fitted at all.
\end{abstract}

\tableofcontents

\section{The problem}

Let $H \in \R^{n\times n}$ be the observed information over the $n = CK + Q$ regression
coefficients. \texttt{CBMRModel.covariance} does two things with it:
\begin{enumerate}
  \item computes $\cond_2(H) = \sigma_{\max}/\sigma_{\min}$ with \texttt{np.linalg.cond}, purely
        to decide whether to warn that the standard errors are untrustworthy;
  \item inverts it with \texttt{np.linalg.inv}.
\end{enumerate}
\texttt{np.linalg.cond} forms a full singular value decomposition. At $n = 7{,}329$ that is
$163$~s, against $14$~s for the inverse it is diagnosing and $1.6$~s for the Hessian itself.
The diagnostic costs an order of magnitude more than the quantity.

\section{Four strategies}

Write $L \in \R^{P\times C}$ for the loadings, $\supp(p) = \{c : L_{pc}\neq 0\}$, and let $\sim$
be the transitive closure of $c \sim c' \iff \exists p:\ c,c'\in\supp(p)$. \begin{proposition}[Structural separation]\label{prop:separation}
For every one of the three distributions, $H_{(ck),(c'k')} = 0$ whenever $c \not\sim c'$.
\end{proposition}
\begin{proof}
The companion note shows each distribution's Hessian has $\beta\beta$ block
$\sum_p L_{pc}L_{pc'}(B^\top\Sigma_pB)_{kk'}$. Every summand carries the factor
$L_{pc}L_{pc'}$, which vanishes unless $c,c' \in \supp(p)$; if some $p$ has both, then
$c \sim c'$.
\end{proof}

That sparsity is the raw material here. It is a statement about the design, not the observation
model, and the entries are structurally absent rather than merely small.

\begin{proposition}[Block-diagonal spectrum]\label{prop:spectrum}
If $H = \blkdiag(H_1,\dots,H_M)$ then $\spec(H) = \bigsqcup_j \spec(H_j)$, and the singular
values are likewise the multiset union. Hence
$\cond_2(H) = \max_j \sigma_{\max}(H_j)\big/\min_j\sigma_{\min}(H_j)$, exactly.
\end{proposition}
\begin{proof}
$H - \lambda I$ is block diagonal with blocks $H_j - \lambda I$, and the determinant of a block
diagonal matrix is the product of its blocks' determinants, so the characteristic polynomial
factorises. Applying the same to $H^\top H$ gives the singular values.
\end{proof}

\begin{proposition}[Symmetry suffices]\label{prop:eigh}
$H$ is symmetric by construction. For symmetric $H$ with eigenvalues $\lambda_i$ the singular
values are $|\lambda_i|$, so $\cond_2(H) = \max_i|\lambda_i|/\min_i|\lambda_i|$ and a symmetric
eigensolver may replace the SVD. No hypothesis on the design, and no positive definiteness.
\end{proposition}
\begin{proof}
Write $H = V\Lambda V^\top$ with $V$ orthogonal; then $H^\top H = V\Lambda^2 V^\top$, so the
singular values --- the nonnegative square roots of the eigenvalues of $H^\top H$ --- are
$|\lambda_i|$.
\end{proof}

\begin{lemma}[Blockwise inverse]\label{lem:blockwise}
If $H = \blkdiag(H_1,\dots,H_M)$ with each $H_j$ nonsingular, then
$H^{-1} = \blkdiag(H_1^{-1},\dots,H_M^{-1})$.
\end{lemma}
\begin{proof}
The product $\blkdiag(H_j)\blkdiag(H_j^{-1})$ is block diagonal with $j$-th block
$H_jH_j^{-1} = I$, so it is the identity; inverses are unique.
\end{proof}

\begin{proposition}[Bordered inverse]\label{prop:schur}
Suppose the spatial columns separate but $Q > 0$, so that
\[
  H = \begin{pmatrix} D & C\\ C^\top & E\end{pmatrix},
  \qquad D = \blkdiag(D_1,\dots,D_M)\in\R^{CK\times CK},\quad C\in\R^{CK\times Q}.
\]
If $D$ and $\Sigma = E - C^\top D^{-1}C$ are nonsingular then
\[
  H^{-1} = \begin{pmatrix}
    D^{-1} + D^{-1}C\Sigma^{-1}C^\top D^{-1} & -D^{-1}C\Sigma^{-1}\\
    -\Sigma^{-1}C^\top D^{-1} & \Sigma^{-1}\end{pmatrix},
\]
at cost $O(\sum_j |D_j|^3) + O(Q^3)$ for the factorisations rather than $O((CK+Q)^3)$.
\end{proposition}
\begin{proof}
Standard block inversion; verify by multiplying out $HH^{-1} = I$ blockwise using
$E - C^\top D^{-1}C = \Sigma$.
\end{proof}

\begin{proposition}[One factorisation for both]\label{prop:chol}
If $H \succ 0$ then a single Cholesky factorisation $H = LL^\top$ serves twice: \texttt{pocon}
returns an estimate of $\cond_1(H)$ in $O(n^2)$ from $L$, and \texttt{potri} returns $H^{-1}$ in
$O(n^3/3)$ from the same $L$. Moreover
\begin{equation}\label{eq:norm-equiv}
  n^{-1}\cond_1(H) \;\le\; \cond_2(H) \;\le\; n\,\cond_1(H).
\end{equation}
\end{proposition}
\begin{proof}
The routines are standard. For \eqref{eq:norm-equiv}, apply
$n^{-1/2}\|M\|_1 \le \|M\|_2 \le n^{1/2}\|M\|_1$ to $M = H$ and $M = H^{-1}$ and multiply.
\end{proof}

\section{What the measurement rejected}

Proposition~\ref{prop:chol} is by far the fastest route and \eqref{eq:norm-equiv} makes it look
safe: a factor of $n$ is four of the sixteen available decimal digits. Measured on the three
structures of Table~\ref{tab:strategies}, the \texttt{pocon} estimate reports $90.09$ against a
true $\cond_2$ of $5.0243$, $2.42\times10^{5}$ against $1657.7$, and $1312.4$ against $4.9796$
--- errors of $18\times$, $146\times$ and $264\times$.

\begin{remark}
Every one of those is inside the bound, which is the point. A bound that admits a $264\times$
error admits behaviour that cannot be shown to a user as a condition number. A fit at
$\cond_2 = 10^{14}$ --- badly conditioned but genuinely invertible --- would be reported near
$10^{16}$ and trip a warning saying its standard errors must not be trusted.
\end{remark}

The resolution is to split the two uses of the factorisation. The Cholesky is kept for the
\emph{inverse}, where it agrees with \texttt{np.linalg.inv} to $1.8\times10^{-15}$ and runs
$2$--$3\times$ faster. It is discarded for the \emph{condition number}, which comes instead from
Proposition~\ref{prop:eigh}: exact, fully general, and still $3$--$4\times$ faster than the SVD.

\begin{table}[htbp]
\centering
\caption{Condition number and inverse on synthetic matrices at the 21-group shape, over the
three structures real designs produce. ``dense'' is the stock path. All inverses agree with the
dense one to $\le 4.6\times10^{-15}$ relative. Rows marked $\star$ are what the dispatch adopts.}
\label{tab:strategies}
\resizebox{\textwidth}{!}{%
\begin{tabular}{llrrrr}
\toprule
structure & strategy & $\cond$ reported & $\cond$ time & inverse time & total\\
\midrule
\multirow{4}{*}{\shortstack[l]{block diagonal\\ \texttt{\textasciitilde\ s(a)}, $7329^2$}}
 & dense (SVD $+$ LU) & $5.0243$ & $157.2$~s & $14.2$~s & $171.4$~s\\
 & \texttt{eigvalsh} $+$ Cholesky & $5.0243$ & $36.5$~s & $7.6$~s & $44.1$~s\\
 & Cholesky $+$ \texttt{pocon} & $90.093$ & $2.9$~s & $5.5$~s & $8.4$~s\\
 & $\star$ blockwise & $5.0243$ & $0.4$~s & $0.2$~s & $\mathbf{0.7}$~s\\
\midrule
\multirow{4}{*}{\shortstack[l]{bordered\\ \texttt{\textasciitilde\ s(a) + x * z}, $7331^2$}}
 & dense (SVD $+$ LU) & $1657.7$ & $169.0$~s & $20.7$~s & $189.6$~s\\
 & $\star$ \texttt{eigvalsh} $+$ Schur & $1657.7$ & $48.5$~s & $4.1$~s & $\mathbf{52.7}$~s\\
 & Cholesky $+$ \texttt{pocon} & $2.42\times10^{5}$ & $4.7$~s & $8.7$~s & $13.3$~s\\
 & blockwise & \multicolumn{4}{c}{\emph{not applicable --- the matrix does not separate}}\\
\midrule
\multirow{3}{*}{\shortstack[l]{fully coupled\\ \texttt{\textasciitilde\ sz(a) + sz(b)}, $7329^2$}}
 & dense (SVD $+$ LU) & $4.9796$ & $161.8$~s & $19.8$~s & $181.7$~s\\
 & $\star$ \texttt{eigvalsh} $+$ Cholesky & $4.9796$ & $52.1$~s & $15.8$~s & $\mathbf{67.8}$~s\\
 & Cholesky $+$ \texttt{pocon} & $1312.4$ & $4.4$~s & $8.6$~s & $13.0$~s\\
\bottomrule
\end{tabular}}
\end{table}

\section{The dispatch, and its generality}

\begin{enumerate}
  \item \textbf{blockwise} --- spatial columns separate, $Q=0$. Condition number by
        Proposition~\ref{prop:spectrum}, inverse blockwise.
  \item \textbf{bordered} --- spatial columns separate, $Q>0$. Inverse by
        Proposition~\ref{prop:schur}; the condition number cannot follow, because the spectrum of
        a bordered matrix is not the union of its parts, so it falls back to
        Proposition~\ref{prop:eigh}.
  \item \textbf{Cholesky} --- no exploitable structure but $H\succ0$. Condition number by
        Proposition~\ref{prop:eigh}, inverse by Cholesky.
  \item \textbf{dense} --- $H\not\succ0$; the Cholesky declines and the general inverse runs.
\end{enumerate}

\begin{proposition}[The dispatch is well defined]\label{prop:welldefined}
On every branch the dispatch returns $H^{-1}$ and the exact $\cond_2(H)$, so the rung taken
affects only cost. Moreover rung~2 reduces to rung~1 when $C = 0$.
\end{proposition}
\begin{proof}
Rung~1 returns $H^{-1}$ by Lemma~\ref{lem:blockwise} and $\cond_2(H)$ by
Proposition~\ref{prop:spectrum}; rung~2 by Proposition~\ref{prop:schur} and
Proposition~\ref{prop:eigh}; rung~3 by uniqueness of the Cholesky inverse and
Proposition~\ref{prop:eigh}; rung~4 by definition. For the reduction, set $C=0$ in
Proposition~\ref{prop:schur}: then $\Sigma = E$, the off-diagonal blocks of the stated inverse
vanish and the leading block is $D^{-1}$, which is Lemma~\ref{lem:blockwise}.
\end{proof}

Every rung therefore reports the \emph{same} condition number, to the last printed digit.
Structural claims are checked against the matrix in hand, never assumed: a design that ought to
separate but whose information matrix has an entry outside the implied blocks is inverted
densely, with a warning.

Table~\ref{tab:ladder} verifies the dispatch over the nine-design, three-distribution suite. All
four rungs fire; every one agrees with the dense inverse to within a tolerance scaled by the
matrix's own conditioning. The important negative result is the bordered rung: a global moderator
destroys block diagonality outright --- one block where there were twenty --- while leaving the
spatial columns separable underneath, so the inverse is recovered even though the classification
that motivated it is lost.

\begin{table}[htbp]
\centering
\caption{The dispatch over every design the CBMR grammar produces, crossed with all three
distributions. ``blocks'' counts the components of $\sim$; ``rung'' is what that buys. Relative
error is against \texttt{np.linalg.inv} of the same matrix, worst over the distributions that
accept the design.}
\label{tab:ladder}
\resizebox{\textwidth}{!}{%
\begin{tabular}{llrlrr}
\toprule
design & formula & blocks & rung & $\cond_2$ & inverse rel.\ error\\
\midrule
cell-means factor & \texttt{\textasciitilde\ s(a)} & $4$ & blockwise & $4.8\times10^{4}$ & $4.0\times10^{-13}$\\
full interaction & \texttt{\textasciitilde\ s(a:b)} & $8$ & blockwise & $4.7\times10^{4}$ & $1.6\times10^{-12}$\\
additive main effects & \texttt{\textasciitilde\ sz(a) + sz(b)} & $1$ & Cholesky & $1.9\times10^{5}$ & $1.2\times10^{-12}$\\
main effects $+$ interaction & \texttt{\textasciitilde\ sz(a) + sz(b) + sz(a:b)} & $1$ & dense$^{\dagger}$ & $4.1\times10^{20}$ & ---\\
one global moderator & \texttt{\textasciitilde\ s(a) + x} & $1$ & bordered & $2.5\times10^{7}$ & $9.0\times10^{-13}$\\
global moderator interaction & \texttt{\textasciitilde\ s(a) + x * z} & $1$ & bordered & $2.6\times10^{7}$ & $8.3\times10^{-13}$\\
voxelwise moderator, continuous & \texttt{\textasciitilde\ s(a) + s(x)} & $1$ & Cholesky & $2.3\times10^{5}$ & $6.4\times10^{-13}$\\
voxelwise moderator, coarse & \texttt{\textasciitilde\ s(a) + s(w)} & $1$ & Cholesky & $6.2\times10^{6}$ & $7.6\times10^{-12}$\\
everything at once & \texttt{\textasciitilde\ sz(a) + sz(b) + s(w) + x * z} & $1$ & Cholesky & $5.1\times10^{7}$ & $3.3\times10^{-11}$\\
\bottomrule
\end{tabular}}

\vspace{2pt}
\raggedright\footnotesize
$^{\dagger}$ Rank deficient at this sample size, so the Cholesky declines and the general inverse
runs. No inverse of this matrix is trustworthy by any route, so no comparison is made; the test
asserts the singularity instead.
\end{table}

\section{Equivalence}

Two independent checks, matching the companion note's structure.

\subsection{Symbolic}

\texttt{scripts/symbolic\_covariance.py} restates each proposition over symbols and closes it
with \texttt{sympy}. Every residual is exactly $0$; nothing numerical is involved. (The script
runs structure, then inverses, then spectra; the table is ordered by proposition instead.)

\begin{center}
\begin{tabular}{llc}
\toprule
claim & residual checked & result\\
\midrule
Prop.~\ref{prop:separation} & cross block of two columns sharing no pattern & $0$\\
\addlinespace
Prop.~\ref{prop:spectrum} & $\det(\blkdiag(A,B)-\lambda I) - \det(A-\lambda I)\det(B-\lambda I)$ & $0$\\
\addlinespace
\multirow{2}{*}{Prop.~\ref{prop:eigh}}
 & $H^\top H - H^2$ for symmetric $H$ & $0$\\
 & $\det(H^2-\lambda^2 I) - \det(H-\lambda I)\det(H+\lambda I)$ & $0$\\
\addlinespace
Lem.~\ref{lem:blockwise} & $\blkdiag(A,B)^{-1} - \blkdiag(A^{-1},B^{-1})$ & $0$\\
\addlinespace
\multirow{4}{*}{Prop.~\ref{prop:schur}}
 & $D\,\mathrm{TL} + C\,\mathrm{BL} - I_m$ & $0$\\
 & $D\,\mathrm{TR} + C\,\mathrm{BR}$ & $0$\\
 & $C^\top \mathrm{TL} + E\,\mathrm{BL}$ & $0$\\
 & $C^\top \mathrm{TR} + E\,\mathrm{BR} - I_q$ & $0$\\
\bottomrule
\end{tabular}
\end{center}

\begin{remark}[Why Proposition~\ref{prop:schur} is proved over matrix symbols]
Filling in a $4\times4$ symbolic $D^{-1}$ inside a $6\times6$ product makes \texttt{simplify}
intractable --- the attempt does not terminate. It also proves less. The identity needs only
that $D$ and $\Sigma$ are invertible and that the shapes conform, so $D^{-1}$ and $\Sigma^{-1}$
are left formal and the sole rewrites used are $DD^{-1} = I$ and, through
$E = \Sigma + C^\top D^{-1}C$, the definition of the Schur complement. The result is a proof for
arbitrary conformable blocks rather than for one filled-in instance, and
Lemma~\ref{lem:blockwise} is what licenses computing that $D^{-1}$ blockwise.
\end{remark}

\subsection{Numerical}

\texttt{test\_cbmr\_analytic.py} applies all four rungs to \emph{the same} matrix rather than to
four separate fits, so nothing about a fit can mask a disagreement: a block diagonal
$48\times48$ and the same matrix with a $2$-column border welded on. Blockwise, Cholesky and
both routes to the condition number match the dense ones to $10^{-10}$; setting the border to
zero returns the Schur complement to the blockwise inverse exactly, which is
Proposition~\ref{prop:welldefined}'s reduction claim in arithmetic.
\texttt{test\_cbmr\_designs.py} then exercises the dispatch on real fits, over the suite of
Table~\ref{tab:ladder}.

\section{Memory}

Speed was measured first because it was the visible failure. Memory is the one that decides
whether the model can be fitted at all.

\subsection{Where it goes}

Three allocations dominate, all $O(n^2)$: the information matrix, the covariance returned to the
caller, and whatever working copy the chosen strategy makes. At the real 21-group shape
$n = 8{,}420$, so each is $0.53$~GB. Everything else --- the sparse foci, the
$(V\times K)$ basis, the L-BFGS history, the per-pattern log-intensity --- is an order of
magnitude smaller.

\begin{remark}[A saving that is not there]
\texttt{np.zeros} is lazily backed, and the blockwise inverse writes only the diagonal blocks
--- $19$~MB of a $400$~MB array at the synthetic shape. It is tempting to conclude that a
block-diagonal covariance is nearly free in resident memory. Measured, it is not: writing the
diagonal blocks moves resident memory by $+0.399$~GB, essentially the whole array, because the
blocks are strided across every row and the fault granularity is coarser than the writes.
Reads genuinely do not fault --- two full passes over the array cost $+0.000$~GB --- but that is
no help, since the array is fully resident by then. Budget $n^2$ words, not the blocks.
\end{remark}

\subsection{Measured, per call}

Table~\ref{tab:memory} reports peak resident memory \emph{during} the \texttt{covariance()} call,
sampled at $100$~Hz from \texttt{/proc/self/statm}. \texttt{ru\_maxrss} is the wrong instrument
here: it is a high-water mark for the whole process, so it reports whatever the setup allocated.

\begin{table}[htbp]
\centering
\caption{Peak resident memory during one \texttt{covariance()} call, one process per row.
``matrix'' is $n^2$ words, the unavoidable floor. The last column is the intermediate stock
\texttt{torch.func.hessian} would require, which is why none of these shapes runs on it.}
\label{tab:memory}
\begin{tabular}{lrlrrrr}
\toprule
shape & $n$ & rung & time & peak $\Delta$ & matrix & stock needs\\
\midrule
diagnosis, 10 groups & $4{,}100$ & blockwise & $1.1$~s & $0.18$~GB & $0.13$~GB & $4.9$~GB\\
task, 10 groups & $4{,}270$ & blockwise & $1.8$~s & $0.19$~GB & $0.14$~GB & $7.7$~GB\\
diagnosis, 21 groups & $7{,}329$ & blockwise & $1.3$~s & $0.48$~GB & $0.40$~GB & $9.7$~GB\\
diagnosis, 21 groups $+$ moderator & $7{,}330$ & bordered & $40.6$~s & $1.28$~GB & $0.40$~GB & $9.7$~GB\\
\bottomrule
\end{tabular}
\end{table}

The blockwise rung costs about what the matrix costs. The bordered rung costs roughly three
times that, and is $30\times$ slower: it pays for a full symmetric eigendecomposition of the
whole matrix (Proposition~\ref{prop:eigh} cannot be localised to the blocks) and for a dense
$D^{-1}$ alongside the result. It is still a large improvement on the dense path, but it is the
expensive rung, and a design that avoids moderators avoids it.

\subsection{Measured, end to end}

Table~\ref{tab:e2e} is the number that answers the practical question. The run is the real
21-group diagnosis pool --- every category with at least $200$ studies, excluding Healthy ---
at $4$~mm with \texttt{spline\_spacing=5}: $20$ groups after pooling, $8{,}178$ experiments,
$8{,}420$ coefficients, $421$ bases, $21{,}789$ analysis voxels. It covers everything a user
does: load, fit, covariance, and twenty one-versus-rest contrasts.

\begin{table}[htbp]
\centering
\caption{Complete 21-group run, real data, cumulative peak \texttt{ru\_maxrss} at each stage.
``chunked'' is the previous \texttt{jacrev} path with the dense SVD and LU; ``closed form'' is
the closed-form Hessian with the dispatch of Section~4. Stock NiMARE does not appear because it
does not complete: \texttt{torch.func.hessian} requires a $28.7$~GB intermediate at this shape.}
\label{tab:e2e}
\begin{tabular}{lrrrr}
\toprule
& \multicolumn{2}{c}{chunked autodiff} & \multicolumn{2}{c}{closed form}\\
\cmidrule(lr){2-3}\cmidrule(lr){4-5}
stage & time & peak & time & peak\\
\midrule
load studyset & $1.8$~s & $0.74$~GB & $1.8$~s & $0.74$~GB\\
fit $+$ covariance & $398.7$~s & $4.40$~GB & $27.4$~s & $2.07$~GB\\
20 one-versus-rest contrasts & $368.8$~s & $5.49$~GB & $14.7$~s & $2.67$~GB\\
\midrule
\textbf{total} & $\mathbf{769}$~s & $\mathbf{5.49}$~GB & $\mathbf{44}$~s & $\mathbf{2.67}$~GB\\
\bottomrule
\end{tabular}
\end{table}

$17\times$ faster and half the memory. The contrast stage improves by more than the memoisation
alone predicts; part of that gap is likely contention on the shared machine and should not be
read as a $25\times$ effect of this change.

\subsection{Will it run on a laptop?}

At $2.67$~GB peak, the complete 21-group analysis fits on a $10$--$12$~GB machine with room to
spare --- roughly a quarter of available RAM, leaving the operating system and everything else
untouched. The $10$-group diagnosis and task models are smaller still. Three caveats:

\begin{itemize}
  \item Memory grows as $n^2 = (CK+Q)^2$, and $C$ is the group count. Doubling the groups
        quadruples the covariance: $40$ groups at $421$ bases gives $n = 16{,}840$ and
        $2.11$~GB per matrix. Extrapolating the measured ratio of peak to matrix size on the
        blockwise rung ($\approx\!1.2$), that is roughly $6$--$8$~GB peak --- still inside a
        $12$~GB machine, but no longer with room to spare. This is an extrapolation, not a
        measurement. It is the ceiling to watch, and \texttt{spline\_spacing} moves it as
        directly as the group count does.
  \item Adding a global moderator moves the fit onto the bordered rung, which measured $2.7\times$
        the memory and $30\times$ the time of the blockwise one.
  \item The covariance is returned dense because \texttt{contrasts.py} indexes it densely. That
        $0.53$~GB is the floor; removing it means changing NiMARE.
\end{itemize}

\section{Identification}

At the real shape there are $8{,}420$ coefficients and $8{,}178$ experiments, which invites the
conclusion that the model has more parameters than data. It does not, and the reason is worth
recording because the reassuring version of the argument is also wrong.

CBMR's likelihood is not over experiment rows. It is over voxel-level foci counts, summarised
per pattern, so there are $P \times V = 435{,}780$ count observations --- $52$ per coefficient.
More to the point, the spatial factor term is block diagonal, so the model is not one problem
with $8{,}420$ parameters but twenty independent problems with $421$ each, and each group's
coefficients are identified by that group's own foci. Measured over the real pool, foci per
coefficient runs from $8.1$ (hearing loss, $221$ studies, $3{,}406$ foci) to $65.4$ (Alzheimer's,
$1{,}345$ studies, $27{,}529$ foci), and every per-group block has condition number between
$3.8\times10^{4}$ and $3.0\times10^{5}$ --- far from the $\epsilon^{-1}\approx4.5\times10^{15}$
at which double precision gives out.

\begin{remark}[The diagnostic does not diagnose this]
Under a canonical link the observed information equals the expected information, so for Poisson
$H$ is a function of the fitted intensity and the group sizes and not of the foci directly. A
group can be far too thin for the asymptotics while still producing a perfectly invertible block.
Fitting three groups of sizes $(40,40,40)$, $(40,40,10)$, $(40,40,4)$ and $(40,40,2)$, the thin
group's condition number stays between $1.3\times10^{4}$ and $1.5\times10^{5}$ throughout ---
no trend, no warning --- while its median coefficient standard error rises from $27.5$ to
$150.2$. The standard errors do respond; the conditioning warning does not. It is the wrong
instrument for detecting a data-starved group, and \texttt{incidence\_threshold} together with
the per-term parameter budget from \texttt{describe\_terms()} is the right one.
\end{remark}

\clearpage
% Collects the machine-generated proof fragments into one appendix.
\section{Generated proof fragments}
\label{sec:generated}

Everything below is written by \texttt{run\_proofs.py} from the same sympy objects the proofs
check. Nothing here is transcribed by hand, so the document cannot disagree with the
verification that backs it.

% Generated by cbmr-proofs. Do not edit.
% Every expression below is rendered from the object sympy verified, not transcribed.

\subsection{Poisson}

The log-likelihood on marginals, dropping the parameter-free
$-\sum\log y!$, is $\log\mathcal{L}=\sum_{v}Y_{v}S_{v}+\sum_i y_i m_i-(\sum_v e^{S_v})T$ with
$T=\sum_i e^{m_i}$. The first two terms are linear in the coefficients, so they differentiate
away entirely; that is the whole of the data dependence.


\noindent The expressions the claims are formed from:

\begin{equation*}
  -\log\mathcal{L} = - Y_{0} \left(B_{00} \beta_{0} + B_{01} \beta_{1}\right) - Y_{1} \left(B_{10} \beta_{0} + B_{11} \beta_{1}\right) - Y_{2} \left(B_{20} \beta_{0} + B_{21} \beta_{1}\right) - y_{0} \left(G_{00} \gamma_{0} + G_{01} \gamma_{1}\right) - y_{1} \left(G_{10} \gamma_{0} + G_{11} \gamma_{1}\right) + \left(e^{G_{00} \gamma_{0} + G_{01} \gamma_{1}} + e^{G_{10} \gamma_{0} + G_{11} \gamma_{1}}\right) \left(e^{B_{00} \beta_{0} + B_{01} \beta_{1}} + e^{B_{10} \beta_{0} + B_{11} \beta_{1}} + e^{B_{20} \beta_{0} + B_{21} \beta_{1}}\right)
\end{equation*}

\noindent Residuals, each asserted to vanish:

\begin{center}\small
\begin{tabular}{lll}
\toprule
claim & residual & result\\
\midrule
d2/dbeta0 dbeta0 & $\partial^2_{\beta_0\beta_0} - T\,(B^\top\mathrm{diag}(u)B)_{00}$ & $0$\\
d2/dbeta0 dbeta1 & $\partial^2_{\beta_0\beta_1} - T\,(B^\top\mathrm{diag}(u)B)_{01}$ & $0$\\
d2/dbeta1 dbeta0 & $\partial^2_{\beta_1\beta_0} - T\,(B^\top\mathrm{diag}(u)B)_{10}$ & $0$\\
d2/dbeta1 dbeta1 & $\partial^2_{\beta_1\beta_1} - T\,(B^\top\mathrm{diag}(u)B)_{11}$ & $0$\\
d2/dbeta0 dgamma0 & $\partial^2_{\beta_0\gamma_0} - a_0U_0$ & $0$\\
d2/dbeta0 dgamma1 & $\partial^2_{\beta_0\gamma_1} - a_0U_1$ & $0$\\
d2/dbeta1 dgamma0 & $\partial^2_{\beta_1\gamma_0} - a_1U_0$ & $0$\\
d2/dbeta1 dgamma1 & $\partial^2_{\beta_1\gamma_1} - a_1U_1$ & $0$\\
d2/dgamma0 dgamma0 & $\partial^2_{\gamma_0\gamma_0} - E\sum_i e_iG_{i0}G_{i0}$ & $0$\\
d2/dgamma0 dgamma1 & $\partial^2_{\gamma_0\gamma_1} - E\sum_i e_iG_{i0}G_{i1}$ & $0$\\
d2/dgamma1 dgamma0 & $\partial^2_{\gamma_1\gamma_0} - E\sum_i e_iG_{i1}G_{i0}$ & $0$\\
d2/dgamma1 dgamma1 & $\partial^2_{\gamma_1\gamma_1} - E\sum_i e_iG_{i1}G_{i1}$ & $0$\\
Hessian contains no foci symbol & $\{Y_v, y_i\} \cap \mathrm{free}(H)$ & $0$\\
\bottomrule
\end{tabular}
\end{center}

\noindent\textbf{Verdict:} every residual is exactly zero.

% Generated by cbmr-proofs. Do not edit.
% Every expression below is rendered from the object sympy verified, not transcribed.

\subsection{Negative Binomial}

NiMARE's negative binomial uses shape $r=s_1^2/(\theta s_2)$ and success
probability $\pi_v=(1+s_1/(\theta u_v s_2))^{-1}$. Under the reparameterisation
$A=s_1/(\theta s_2)$, $R=s_1A$ we have $\pi_v=u_v/(u_v+A)$, and the special functions depend on
$R$ alone, the voxels on $S$ alone, and the moderators on $(R,A)$ alone.


\noindent The expressions the claims are formed from:

\begin{equation*}
  \Psi(R, A, S) = - 3 R \log{\left(A \right)} - S_{0} y_{0} - S_{1} y_{1} - S_{2} y_{2} + \left(R + y_{0}\right) \log{\left(A + e^{S_{0}} \right)} + \left(R + y_{1}\right) \log{\left(A + e^{S_{1}} \right)} + \left(R + y_{2}\right) \log{\left(A + e^{S_{2}} \right)} + 3 \operatorname{loggamma}{\left(R \right)} - \operatorname{loggamma}{\left(R + y_{0} \right)} - \operatorname{loggamma}{\left(R + y_{1} \right)} - \operatorname{loggamma}{\left(R + y_{2} \right)}
\end{equation*}

\noindent Residuals, each asserted to vanish:

\begin{center}\small
\begin{tabular}{lll}
\toprule
claim & residual & result\\
\midrule
Psi(R,A,S) equals NiMARE's summand & $\Psi\big|_{R,A} - \big(-\log\mathcal{L} - \textrm{const}\big)$ & $0$\\
d2Psi/dS0\^{}2 & $\partial^2_{S_0}\Psi - \frac{(R+Y_0)u_0A}{(u_0+A)^2}$ & $0$\\
d2Psi/dS0 dR & $\partial^2_{S_0R}\Psi - \frac{u_0}{u_0+A}$ & $0$\\
d2Psi/dS0 dA & $\partial^2_{S_0A}\Psi + \frac{(R+Y_0)u_0}{(u_0+A)^2}$ & $0$\\
d2Psi/dS1\^{}2 & $\partial^2_{S_1}\Psi - \frac{(R+Y_1)u_1A}{(u_1+A)^2}$ & $0$\\
d2Psi/dS1 dR & $\partial^2_{S_1R}\Psi - \frac{u_1}{u_1+A}$ & $0$\\
d2Psi/dS1 dA & $\partial^2_{S_1A}\Psi + \frac{(R+Y_1)u_1}{(u_1+A)^2}$ & $0$\\
d2Psi/dS2\^{}2 & $\partial^2_{S_2}\Psi - \frac{(R+Y_2)u_2A}{(u_2+A)^2}$ & $0$\\
d2Psi/dS2 dR & $\partial^2_{S_2R}\Psi - \frac{u_2}{u_2+A}$ & $0$\\
d2Psi/dS2 dA & $\partial^2_{S_2A}\Psi + \frac{(R+Y_2)u_2}{(u_2+A)^2}$ & $0$\\
Psi\_R & $\Psi_R\ \textrm{as implemented}$ & $0$\\
Psi\_A & $\Psi_A\ \textrm{as implemented}$ & $0$\\
Psi\_RR & $\Psi_{RR}\ \textrm{as implemented}$ & $0$\\
Psi\_RA & $\Psi_{RA}\ \textrm{as implemented}$ & $0$\\
Psi\_AA & $\Psi_{AA}\ \textrm{as implemented}$ & $0$\\
dR/dgamma0 & $\partial_{\gamma_0}R - R(2\ell_0-h_0)$ & $0$\\
dA/dgamma0 & $\partial_{\gamma_0}A - A(\ell_0-h_0)$ & $0$\\
d2R/dgamma0 dgamma0 & $\partial^2_{\gamma_0\gamma_0}R\ \textrm{as implemented}$ & $0$\\
d2A/dgamma0 dgamma0 & $\partial^2_{\gamma_0\gamma_0}A\ \textrm{as implemented}$ & $0$\\
d2R/dgamma0 dgamma1 & $\partial^2_{\gamma_0\gamma_1}R\ \textrm{as implemented}$ & $0$\\
d2A/dgamma0 dgamma1 & $\partial^2_{\gamma_0\gamma_1}A\ \textrm{as implemented}$ & $0$\\
dR/dgamma1 & $\partial_{\gamma_1}R - R(2\ell_1-h_1)$ & $0$\\
dA/dgamma1 & $\partial_{\gamma_1}A - A(\ell_1-h_1)$ & $0$\\
d2R/dgamma1 dgamma0 & $\partial^2_{\gamma_1\gamma_0}R\ \textrm{as implemented}$ & $0$\\
d2A/dgamma1 dgamma0 & $\partial^2_{\gamma_1\gamma_0}A\ \textrm{as implemented}$ & $0$\\
d2R/dgamma1 dgamma1 & $\partial^2_{\gamma_1\gamma_1}R\ \textrm{as implemented}$ & $0$\\
d2A/dgamma1 dgamma1 & $\partial^2_{\gamma_1\gamma_1}A\ \textrm{as implemented}$ & $0$\\
\bottomrule
\end{tabular}
\end{center}

\noindent\textbf{Verdict:} every residual is exactly zero.

% Generated by cbmr-proofs. Do not edit.
% Every expression below is rendered from the object sympy verified, not transcribed.

\subsection{Clustered Negative Binomial}

The latent effect is a property of the experiment rather than of the voxel, so
the spatial coefficients reach the likelihood only through the scalar $E=\sum_v u_v$. Writing
$\Gamma(E)=\sum_i (y_i+\nu)\log(Ee_i+\nu)$, the second differential in $S$ is
$\Gamma''(E)u_vu_{v'}+\Gamma'(E)u_v\delta_{vv'}$ --- diagonal plus rank one.


\noindent The expressions the claims are formed from:

\begin{equation*}
  -\log\mathcal{L}_p = - Y_{0} \left(B_{00} \beta_{0} + B_{01} \beta_{1}\right) - Y_{1} \left(B_{10} \beta_{0} + B_{11} \beta_{1}\right) - Y_{2} \left(B_{20} \beta_{0} + B_{21} \beta_{1}\right) - 2 \nu \log{\left(\nu \right)} - y_{0} \left(G_{00} \gamma_{0} + G_{01} \gamma_{1}\right) - y_{1} \left(G_{10} \gamma_{0} + G_{11} \gamma_{1}\right) + \left(\nu + y_{0}\right) \log{\left(\nu + \left(e^{B_{00} \beta_{0} + B_{01} \beta_{1}} + e^{B_{10} \beta_{0} + B_{11} \beta_{1}} + e^{B_{20} \beta_{0} + B_{21} \beta_{1}}\right) e^{G_{00} \gamma_{0} + G_{01} \gamma_{1}} \right)} + \left(\nu + y_{1}\right) \log{\left(\nu + \left(e^{B_{00} \beta_{0} + B_{01} \beta_{1}} + e^{B_{10} \beta_{0} + B_{11} \beta_{1}} + e^{B_{20} \beta_{0} + B_{21} \beta_{1}}\right) e^{G_{10} \gamma_{0} + G_{11} \gamma_{1}} \right)} + 2 \operatorname{loggamma}{\left(\nu \right)} - \operatorname{loggamma}{\left(\nu + y_{0} \right)} - \operatorname{loggamma}{\left(\nu + y_{1} \right)}
\end{equation*}
\begin{equation*}
  g' = \frac{\left(\nu + y_{0}\right) e^{G_{00} \gamma_{0} + G_{01} \gamma_{1}}}{\nu + \left(e^{B_{00} \beta_{0} + B_{01} \beta_{1}} + e^{B_{10} \beta_{0} + B_{11} \beta_{1}} + e^{B_{20} \beta_{0} + B_{21} \beta_{1}}\right) e^{G_{00} \gamma_{0} + G_{01} \gamma_{1}}} + \frac{\left(\nu + y_{1}\right) e^{G_{10} \gamma_{0} + G_{11} \gamma_{1}}}{\nu + \left(e^{B_{00} \beta_{0} + B_{01} \beta_{1}} + e^{B_{10} \beta_{0} + B_{11} \beta_{1}} + e^{B_{20} \beta_{0} + B_{21} \beta_{1}}\right) e^{G_{10} \gamma_{0} + G_{11} \gamma_{1}}}
\end{equation*}
\begin{equation*}
  g'' = - \frac{\left(\nu + y_{0}\right) e^{2 G_{00} \gamma_{0} + 2 G_{01} \gamma_{1}}}{\left(\nu + \left(e^{B_{00} \beta_{0} + B_{01} \beta_{1}} + e^{B_{10} \beta_{0} + B_{11} \beta_{1}} + e^{B_{20} \beta_{0} + B_{21} \beta_{1}}\right) e^{G_{00} \gamma_{0} + G_{01} \gamma_{1}}\right)^{2}} - \frac{\left(\nu + y_{1}\right) e^{2 G_{10} \gamma_{0} + 2 G_{11} \gamma_{1}}}{\left(\nu + \left(e^{B_{00} \beta_{0} + B_{01} \beta_{1}} + e^{B_{10} \beta_{0} + B_{11} \beta_{1}} + e^{B_{20} \beta_{0} + B_{21} \beta_{1}}\right) e^{G_{10} \gamma_{0} + G_{11} \gamma_{1}}\right)^{2}}
\end{equation*}

\noindent Residuals, each asserted to vanish:

\begin{center}\small
\begin{tabular}{lll}
\toprule
claim & residual & result\\
\midrule
d2/dbeta0 dbeta0 & $\partial^2_{\beta_0\beta_0} - \textrm{closed form}$ & $0$\\
d2/dbeta0 dbeta1 & $\partial^2_{\beta_0\beta_1} - \textrm{closed form}$ & $0$\\
d2/dbeta0 dgamma0 & $\partial^2_{\beta_0\gamma_0} - \textrm{closed form}$ & $0$\\
d2/dbeta0 dgamma1 & $\partial^2_{\beta_0\gamma_1} - \textrm{closed form}$ & $0$\\
d2/dbeta1 dbeta1 & $\partial^2_{\beta_1\beta_1} - \textrm{closed form}$ & $0$\\
d2/dbeta1 dgamma0 & $\partial^2_{\beta_1\gamma_0} - \textrm{closed form}$ & $0$\\
d2/dbeta1 dgamma1 & $\partial^2_{\beta_1\gamma_1} - \textrm{closed form}$ & $0$\\
d2/dgamma0 dgamma0 & $\partial^2_{\gamma_0\gamma_0} - \textrm{closed form}$ & $0$\\
d2/dgamma0 dgamma1 & $\partial^2_{\gamma_0\gamma_1} - \textrm{closed form}$ & $0$\\
d2/dgamma1 dgamma1 & $\partial^2_{\gamma_1\gamma_1} - \textrm{closed form}$ & $0$\\
\bottomrule
\end{tabular}
\end{center}

\noindent\textbf{Verdict:} every residual is exactly zero.

% Generated by cbmr-proofs. Do not edit.
% Every expression below is rendered from the object sympy verified, not transcribed.

\subsection{Covariance strategies}

Let $\mathrm{supp}(p)=\{c: L_{pc}\neq0\}$ and let $\sim$ be the transitive
closure of $c\sim c' \iff \exists p:\,c,c'\in\mathrm{supp}(p)$. Every distribution's Hessian has
$\beta\beta$ block $\sum_p L_{pc}L_{pc'}(B^\top\Sigma_pB)_{kk'}$, so the factor $L_{pc}L_{pc'}$
decides which entries can be nonzero at all.


\noindent Residuals, each asserted to vanish:

\begin{center}\small
\begin{tabular}{lll}
\toprule
claim & residual & result\\
\midrule
columns sharing no pattern have a zero cross block & $H_{(0k),(2k')}\ \textrm{with}\ 0\not\sim 2$ & $0$\\
characteristic polynomial factorises over the blocks & $\det(\mathrm{blkdiag}(A,B)-\lambda I)-\det(A-\lambda I)\det(B-\lambda I)$ & $0$\\
symmetric H satisfies H\^{}T H = H\^{}2 & $H^\top H - H^2$ & $0$\\
the singular values are the absolute eigenvalues & $\det(H^2-\lambda^2I)-\det(H-\lambda I)\det(H+\lambda I)$ & $0$\\
blockwise inverse & $\mathrm{blkdiag}(A,B)^{-1}-\mathrm{blkdiag}(A^{-1},B^{-1})$ & $0$\\
D.TL + C.BL = I & $D\,\mathrm{TL}+C\,\mathrm{BL}-I_m$ & $0$\\
D.TR + C.BR = 0 & $D\,\mathrm{TR}+C\,\mathrm{BR}$ & $0$\\
C\^{}T.TL + E.BL = 0 & $C^\top\mathrm{TL}+E\,\mathrm{BL}$ & $0$\\
C\^{}T.TR + E.BR = I & $C^\top\mathrm{TR}+E\,\mathrm{BR}-I_q$ & $0$\\
\bottomrule
\end{tabular}
\end{center}

\noindent\textbf{Verdict:} every residual is exactly zero.



\end{document}
